\documentclass[12pt]{exam}

\usepackage[margin=1in]{geometry}

\newcommand{\class}{PHYS 110}
\newcommand{\term}{Fall 2023}
\newcommand{\examnum}{C10 (v2)}
\newcommand{\examdate}{\today}
\newcommand{\timelimit}{60 Minutes}


% Engine-specific settings
% Detect pdftex/xetex/luatex, and load appropriate font packages.
% This is inspired by the approach in the iftex package.
% pdftex:
\ifx\pdfmatch\undefined
\else
    \usepackage[T1]{fontenc}
    \usepackage[utf8]{inputenc}
\fi
% xetex:
\ifx\XeTeXinterchartoks\undefined
\else
    \usepackage{fontspec}
    \defaultfontfeatures{Ligatures=TeX}
\fi
% luatex:
\ifx\directlua\undefined
\else
    \usepackage{fontspec}
\fi
% End engine-specific settings

\usepackage{amsmath,amssymb}
\usepackage{fullpage}
\usepackage{graphicx}
\usepackage[svgnames]{xcolor}
\usepackage{url}


\usepackage{tikz}
\usepackage{pgfplots}
\pgfplotsset{compat=1.14}

\urlstyle{same}

\usepackage{pythontex}
% \restartpythontexsession{\thesection}


\usepackage[framemethod=TikZ]{mdframed}

\newcommand{\pytex}{Python\TeX}
\renewcommand*{\thefootnote}{\fnsymbol{footnote}}

%\usepackage{circuitikz}
\usepackage{siunitx}
%\usepackage{pgfplots}
\usepackage{nopageno}


%\pgfplotsset{compat=1.17}
\begin{document}

\pgfplotsset{
  axis lines=middle,
  axis line style={draw=none},
  xmin=0,
  xmax=4.5*pi,
  xticklabels=empty,
  xtick=\empty,
  ymin=-0.25,
  ymax=0.25,
  yticklabels=empty,
  ytick=\empty,
  samples=250,
  domain=0:4.5*pi,
  every axis plot post/.append style={very thick},
  clip=false, width=14cm,height=6cm
             }% end of common axis set



\begin{pycode}

from pylab import *
import random
from randassign import RandAssign
ra = RandAssign()

def make_problem():
	# Random values for this problem
	def rand_tension():
		return round(uniform(1.9,8.1),2)
	def rand_length():
		return round(uniform(0.65,1.25),2)
	def rand_mu():
		return round(round(uniform(100,1900),2)/1000000,5)
	def rand_number():
		return random.randrange(3,6)	

		
	tension = rand_tension()
	length = rand_length()
	mu = rand_mu()
	number = rand_number()


	## Put the problem text here:	
	problem_template = r'''
	The diagram below shows a standing wave on a string which is fixed at both ends.
	\begin{{tikzpicture}}

	\begin{{axis}}% 1. plot

	\addplot[red] {{(.1255*sin(deg({3}*x/4.5)))}};
	\addplot[red,dashed] {{(-.125*sin(deg({3}*x/4.5)))}};
	\draw[thick,<->,blue] (0,-.15) -- (4.5*pi,-.15) node[left=85,below=7] {{ L = $\SI{{{0:.3g}}}{{\meter}}$}};

	\addplot [only marks, gray]
        coordinates {{ (0,0) (4.5*pi,0) }};

	\end{{axis}}
	\end{{tikzpicture}}
	
	The length of the string is $\SI{{{0:.3g}}}{{\meter}}$. 
	The string has a linear mass density of $\mu=\SI{{{1:.3g}}}{{\kilogram/\meter}}$.
	The tension in the string is $\SI{{{2:.3g}}}{{\newton}}$.
	Find the frequencies and wavelengths for the first ${{{3:.3g}}}$ vibration modes of the string.
	'''
	
	## SOLUTION GOES HERE
	c = sqrt(tension/mu)
	i=1
	freq = []
	while i<number:
		freq.append(i*c/(2*length))
		i = i+1
	


	problem = problem_template.format(length,mu,tension,number)
	
	return(problem,freq)
#	return(problem)
	
## Generate 4 unique problems
## The chance of two identical problems is very low in this case
## But checking for identical problems won't hurt
problems = []
while len(problems) < 4:
    p,f = make_problem()
    if p not in problems:
        problems.append(p)
        t = '\\SI{{{0:.3g}}}{{\\hertz}}'
        #['{:.2f}'.format(i) for i in f]
        ra.addsoln(['\\SI{{{0:.3g}}}{{\\hertz}}'.format(i) for i in f])

print('\\noindent \\begin{tabular}{l l r} \\textbf{\\class} - \\textbf{\\examnum}  & \\textbf{Name} \\makebox[3.25in]{\\hrulefill}  \\\\ \\end{tabular} \\\\ \\rule[1ex]{\\textwidth}{2.21pt}')

print('\\begin{enumerate}')
for n, p in enumerate(problems):
    #print('\\noindent \\begin{tabular}{l l r} \\textbf{\\class} - \\textbf{\\examnum}  & \\textbf{Name} \\makebox[3.25in]{\\hrulefill}  \\\\ \\end{tabular} \\\\ \\rule[1ex]{\\textwidth}{2.21pt} \\begin{questions}\\question '+ p + '\\end{questions}\\newpage')
    #print('\\begin{questions}\\question '+ p + '\\end{questions}\\newpage')
    print('\\item ' + p + '\\newpage')
    if n+1 < len(problems):
        print(r'\vspace{1in}')
print('\\end{enumerate}')


\end{pycode}


	
\end{document}
